GPU performance engineering is expert work. Getting a kernel to run faster requires reasoning about the
memory hierarchy, warp scheduling, occupancy and instruction-level behaviour at the same time, and
expert tuners routinely spend weeks on a single kernel. Large language models can attempt the same edit
in seconds, which raises an obvious question: how much of that expertise can be supplied through the
prompt, and how much of the resulting speedup survives contact with a real application?

Existing code benchmarks do not answer this. HumanEval-style suites \cite{chen2021codex} and
repository-level issue resolution \cite{jimenez2024swebench} score functional correctness, not
execution time, and recent GPU-specific benchmarks \cite{ouyang2025kernelbench,nichols2024parallel}
evaluate standalone kernels written from a reference implementation rather than edits to code that a
scientific application already depends on. Meanwhile, the first thing a human expert looks at---profiler
output---is exactly the kind of context an LLM never sees by default.

This paper studies an end-to-end loop that closes that gap: profile the application, summarize the
profile into the prompt, ask the model for an optimized version of one source file, build it, verify it
against the application's own notion of correctness, measure it, and feed failures back for a bounded
number of repair iterations. Within this loop we treat the prompt as the object of study. Our
contributions are:

\begin{itemize}\itemsep2pt
\item \textbf{A prompt-component ablation} over six variants, two models and six HeCBench kernels
  (Section~\ref{sec:ablation}--\ref{sec:results}), showing that prompt components are not equally
  important: stripping context (\texttt{minimal}) consistently degrades results, while an explicit
  reasoning step (\texttt{cot}) is the only variant that never hurts.
\item \textbf{A semantic-preservation constraint} that suppresses benchmark gaming
  (Section~\ref{sec:intent}). Without it, both models delete the very computation a benchmark exists to
  measure and report speedups of $1504\times$ and $2809\times$ while output and the built-in
  verification remain unchanged; with it, both return honest results. We also report the detection rule
  we use, since neither output comparison nor the benchmark's own check can catch this class of edit.
\item \textbf{An end-to-end evaluation on three production applications}---AutoDock-GPU, LAMMPS and
  Cholla (Section~\ref{sec:realapps})---with per-application correctness gates that we validate by
  fault injection, and interleaved measurement on shared GPUs. The best model output reaches
  $1.571\times$ on Cholla, beating our expert hand-tuned reference, while on AutoDock-GPU no model or
  prompt found the restructuring that is worth $1.532\times$ by hand.
\item \textbf{A profile-data ablation} (Section~\ref{sec:profabl}) that isolates the contribution of
  PC-sampling and roofline data by removing it from an otherwise identical prompt. Profiling data is
  worth $7.7$--$11.5\%$ on the large application and nothing measurable on single-kernel benchmarks,
  which suggests its value scales with how hard the bottleneck is to locate from source alone.
\item \textbf{A measurement of what happens when profiling data is made prescriptive}
  (Section~\ref{sec:guidvar}). Converting metrics into a checklist of named remedies produces both our
  largest gain ($1.04\times \to 1.38\times$) and our largest loss ($7.84\times \to 1.08\times$, where the
  checklist displaced a restructuring the model had found unaided), at repetition spreads below $5\%$.
  Interpretation widens the distribution of outcomes rather than shifting its centre, which a mean
  speedup hides.
\item \textbf{A set of failure modes specific to real code} (Section~\ref{sec:defects}), including compilation-context blindness
  (the model cannot see that its file is textually included elsewhere, or that ten call sites depend on
  a signature), applications whose error handling discards the diagnostic the repair loop needs, and
  measurement pitfalls on shared machines that silently invalidate results.
\end{itemize}

Taken together, the results argue that prompt design is a first-order lever for turning general-purpose
code models into usable performance tools---but also that conclusions drawn on single-kernel benchmarks
do not automatically transfer to the applications those kernels are meant to represent.
